\section{Experimental Setup}
\label{sec:experimental_setup}

\subsection{Datasets}
\label{sec:datasets}

We draw evaluation instances from two benchmark families:

\paragraph{SWE-bench} \citep{swebench2024} provides real GitHub issues paired with developer-written patches. Each instance specifies a repository, a \texttt{base\_commit} at which the issue exists, and a gold patch from which we extract gold file labels. The full dataset contains 2{,}294 test instances, 225 dev instances, and 19{,}008 training instances across 12 repositories. We use the test split for evaluation.

\paragraph{SWE-PolyBench} \citep{swepolybench2025} extends repository-level evaluation to multiple programming languages (Java, JavaScript, TypeScript, Python) and provides both retrieval-metric and execution-based evaluation modes. The Verified subset contains 382 instances. We use the Python subset for cross-benchmark validation.

\subsection{Repository Selection}
\label{sec:repo_selection}

We evaluate on seven repositories spanning four application domains, drawn from both SWE-bench and SWE-PolyBench:

\begin{table*}[t]
\centering
\small
\begin{tabular}{llllr}
\toprule
Repository & Domain & Benchmark & Source prefix & $n$ \\
\midrule
\texttt{pallets/flask} & Web framework & SWE-bench & \texttt{src/flask} & 10 \\
\texttt{psf/requests} & Networking & SWE-bench & \texttt{requests} & 10 \\
\texttt{pytest-dev/pytest} & Testing/tooling & SWE-bench & \texttt{src/\_pytest} & 10 \\
\texttt{tiangolo/fastapi} & Service framework & SWE-bench & \texttt{fastapi} & 10 \\
\texttt{yt-dlp/yt-dlp} & App/CLI & SWE-PolyBench & \texttt{yt\_dlp} & 5 \\
\texttt{langchain-ai/langchain} & Service framework & SWE-PolyBench & \texttt{libs/langchain} & 10 \\
\texttt{keras-team/keras} & Library & SWE-PolyBench & \texttt{keras} & 10 \\
\bottomrule
\end{tabular}
\caption{Repositories in the evaluation matrix. $n$ = number of issues per repository per track. Source prefix restricts graph construction and retrieval to relevant subdirectories.}
\label{tab:repos}
\end{table*}

The selection covers both small, well-structured libraries (Flask, Requests) and larger, more complex codebases (Keras, LangChain), providing diversity in graph structure and retrieval difficulty.

\subsection{Issue Sampling and Manifest Pinning}
\label{sec:issue_sampling}

For each repository, we select issues using a deterministic protocol:
\begin{enumerate}[leftmargin=1.5em]
  \item Sweep all available splits (test, dev, train) and collect instances matching the target repository.
  \item Deduplicate by \texttt{instance\_id}.
  \item Select the first $n$ instances in dataset order.
  \item Pin the resulting instance IDs in a YAML manifest with a unique \texttt{issue\_set\_id}.
\end{enumerate}

The manifest is version-controlled and used for all subsequent runs, ensuring that different methods and repeat runs are evaluated on exactly the same issue set. Example: \texttt{issues\_swebench\_requests\_v2\_10}.

\subsection{Evaluation Tracks}
\label{sec:eval_tracks}

We report results on two parallel tracks to separate retrieval quality from amortization effects:

\paragraph{Track 1: Strict commit fidelity.}
Each issue is evaluated at its own \texttt{base\_commit}. The repository is checked out and the graph/index is rebuilt for each unique commit. This is the ``realistic conservative'' setting: it measures retrieval quality without any amortization benefit. Issues are grouped by commit to avoid redundant rebuilds when multiple issues share a commit, but in practice the commit repeat ratio is low (often 0\%) for $n=10$ samples.

\paragraph{Track 2: Same-snapshot amortized.}
All issues for a given repository are evaluated against a single fixed snapshot (a pinned commit hash). The graph and index are built once and reused for all issues. This measures the ``world-model reuse'' scenario where a pre-built structural representation serves multiple retrieval tasks. The \texttt{commit\_repeat\_ratio} and \texttt{cache\_hit\_rate} are reported to quantify the amortization opportunity.

Results from the two tracks are always reported separately and never merged into a single aggregate.

\subsection{Metrics}
\label{sec:metrics}

\paragraph{Retrieval quality.}
We compute file-level precision, recall, and F1 per issue:
\begin{equation}
P_i = \frac{|\hat{F}_i \cap F_i^*|}{|\hat{F}_i|}, \quad
R_i = \frac{|\hat{F}_i \cap F_i^*|}{|F_i^*|}, \quad
F1_i = \frac{2 P_i R_i}{P_i + R_i}
\end{equation}
where $\hat{F}_i$ is the set of files retrieved by the method and $F_i^*$ is the gold set extracted from patch diffs. Issues that produce errors (API failures, timeouts) receive $F1 = 0$ to avoid survivorship bias. We report mean values across issues per method.

\paragraph{Cost metrics.}
For each method, we report:
\begin{itemize}[leftmargin=1.5em]
  \item \textbf{Setup embedding tokens} ($C^{\text{setup}}$): embedding tokens consumed during index construction.
  \item \textbf{Runtime LLM tokens} ($C^{\text{llm}}$): language model tokens consumed during per-issue retrieval.
  \item \textbf{Query embedding tokens} ($C^{\text{qemb}}$): embedding tokens consumed during per-issue retrieval.
  \item \textbf{Total cost tokens}: $C^{\text{setup}} + \sum_i (C_i^{\text{llm}} + C_i^{\text{qemb}})$.
\end{itemize}

\paragraph{Amortization descriptors.}
\begin{itemize}[leftmargin=1.5em]
  \item $\texttt{commit\_repeat\_ratio} = 1 - \frac{\text{unique base commits}}{n}$
  \item $\texttt{cache\_hit\_rate} = \frac{\text{reused index builds}}{n}$
\end{itemize}

\subsection{Statistical Protocol}
\label{sec:stats}

To account for LLM stochasticity, each experiment cell is repeated at least 3 times with identical configurations.

\paragraph{Paired analysis.}
For each issue $i$, we compute per-issue deltas between method pairs:
\[
\delta_i = F1_i^{(\text{method A})} - F1_i^{(\text{method B})}
\]
We report the mean delta with a bootstrap 95\% confidence interval (10{,}000 resamples). A comparison is considered significant when the CI excludes zero.

\paragraph{Readiness gates.}
Results are promoted to the manuscript only when:
\begin{enumerate}[leftmargin=1.5em]
  \item At least 3 repeat runs are completed.
  \item Bootstrap CIs are computed for all primary comparisons.
  \item The error rate is below 20\% across methods.
\end{enumerate}

\subsection{Infrastructure}
\label{sec:infrastructure}

\begin{itemize}[leftmargin=1.5em]
  \item \textbf{Retrieval language model}: \texttt{gemini-3-flash-preview} for all LLM-guided retrieval agents.
  \item \textbf{Embedding model}: \texttt{gemini-embedding-001} for all vector embeddings.
  \item \textbf{Vector index}: FAISS \citep{johnson2021faiss} flat inner-product index.
  \item \textbf{AST parsing}: \texttt{tree-sitter} \citep{brunsfeld2018treesitter} with the Python grammar.
  \item \textbf{Graph library}: NetworkX \citep{hagberg2008networkx}.
  \item \textbf{Agent turn budget}: up to 8 turns per issue for both graph and RAG agents.
\end{itemize}

\subsection{Reproducibility}
\label{sec:reproducibility}

All experiment configurations are specified in versioned YAML manifests (e.g., \texttt{experiments\_matrix\_v2.yaml}) with pinned instance IDs and snapshot commits. Run outputs are archived with timestamps, and report-ready numbers are sourced from frozen artifact bundles containing raw results, aggregated summaries, and manifest metadata. The full evaluation pipeline is executable via
\texttt{python run\_suite.py experiments\_matrix\_v2.yaml}.

\subsection{Patching Pilot Setup}
\label{sec:patching_setup}

To validate that retrieval quality translates to end-to-end issue resolution, we conduct a patching pilot using the SWE-bench Verified subset. We draw a frozen split of $N=100$ instances across 9 repositories (astropy, matplotlib, Flask, requests, xarray, pytest, scikit-learn, sphinx-doc, sympy), stratified by per-repository Verified instance count. Instance IDs are pinned in a versioned manifest prior to any runs (\texttt{patch\_manifests/verified\_n100\_split\_v1.json}).

We evaluate four methods:
\begin{itemize}[leftmargin=1.5em]
  \item \textbf{Oracle}: gold patch files fed directly to the patch agent (retrieval ceiling).
  \item \textbf{GM-progressive}: graph-guided retrieval feeds the retrieved file set to the patch agent.
  \item \textbf{RAG-progressive}: RAG retrieval feeds the retrieved file set to the patch agent.
  \item \textbf{None} (cold-start): the patch agent receives only the issue text with no retrieval context.
\end{itemize}

The patch agent uses \texttt{gemini-3-flash-preview} with a single generation attempt per instance (\texttt{patch\_max\_turns=1}), a 65,536-token output budget, a 200K-character file context limit, and a 480-second wall-clock cap per instance (\texttt{instance\_wall\_clock\_cap\_s=480}). Patches are evaluated using the official SWE-bench Docker harness for execution-based pass/fail assessment. Results are explicitly labeled as a pilot study: single run per method, no repeats.

\paragraph{Cost accounting disclosure.}
Frozen GM and RAG runs were produced under a dual-build pipeline in which both graph and RAG indices were built for GM and RAG methods respectively. We therefore report two cost views: \emph{method-accounted} (design intent: GM pays only for graph build; RAG pays only for RAG index build) and \emph{as-run-consumed} (actual API spend including the dual-build overhead). A split-build fix is implemented in \texttt{run\_patch.py} for future runs; the as-run view captures the actual cost incurred during this pilot.

\paragraph{Timeout accounting.}
Instances that exceed the 480-second wall-clock cap are included in resolved-rate denominators (all-in primary metric). A timeout-excluded sensitivity analysis is reported alongside to separate method quality from infrastructure-imposed censoring.
